\section{Conclusion}
\label{sec:conclusion}

\subsection{What we built}
This report walked through a complete three-stage pipeline that takes
\modelname{Qwen2.5-1.5B} from a base model that essentially fails on math
benchmarks under our instruction template to a 1.5B-parameter model that
solves $\sim70\%$ of \code{GSM8K} and $\sim45\%$ of \code{MATH-500}. The
pipeline has three replaceable parts: (i) a small SFT step on
\textbf{1\,298} cleaned reasoning traces, (ii) a verifiable-reward RL
step on up to \textbf{81\,463} math problems using GRPO or DAPO, and
(iii) a composite reward that mixes correctness, format shaping, and
several length-and-cleanliness penalties.

\subsection{Headline takeaways}
\begin{enumerate}
  \item \emph{Most of the math gain comes from RL.} SFT alone explains
    only the first $\sim10\!\times$ improvement on \code{GSM8K} and
    \code{MATH-500}. The much larger second jump comes from the
    verifiable-reward RL stage.
  \item \emph{GRPO and DAPO are nearly equivalent} under our reward
    shape and hyperparameter range. The numerical differences in
    \cref{tab:main} are smaller than the difference between the two
    different training pools.
  \item \emph{General ability is preserved.} Both \code{HellaSwag} and
    \code{MMLU} stay at the base level (\code{HellaSwag} drifts slightly
    upwards, \code{MMLU} stays statistically flat).
  \item \emph{Format failures disappear, semantic ones remain.} After
    RL, the inspected samples no longer show role-token leakage or
    truncated answer lines. The remaining errors are dominated by the
    four semantic patterns in \cref{tab:failures}.
\end{enumerate}

\subsection{Limitations}
\begin{itemize}
  \item \emph{Scale.} A single 1.5B base model on a single GPU does not
    establish how the same recipe transfers to 7B+ regimes. Earlier
    internal experiments at 0.5B already showed that very small bases are
    fragile in both SFT and RL, suggesting non-linear interaction with
    scale.
  \item \emph{Reward shape.} The reward function in \cref{sec:reward} is
    hand-tuned. A more principled comparison against simpler rewards
    (\eg pure $\{0,1\}$ correctness or correctness $+$ format only)
    would clarify which terms are necessary versus convenient.
  \item \emph{Checkpoints.} The reported \code{dapo} and \code{grpo}
    rows correspond to step-150 checkpoints rather than the final
    training step, since the reward and length curves had already
    plateaued. A more careful early-stopping protocol could shift the
    headline numbers slightly.
  \item \emph{Failure analysis.} The four-way taxonomy in
    \cref{tab:failures} is qualitative; a quantitative count over a
    larger inspected pool would make the priorities for future reward
    additions sharper.
\end{itemize}

\subsection{Future work}
Two concrete extensions follow naturally from
\cref{sec:failure}: a precondition-aware reward to attack
\emph{overconfident heuristic shortcuts} on \code{TheoremQA}, and a
state-consistency micro-reward to attack \emph{event-boundary} errors on
\code{GSM8K}. Both require an external verifier richer than the current
final-answer check, but the current pipeline cleanly accommodates an
extra reward component without changing the rest of the stack.

\paragraph{Reproducibility.}
All code, prompts, hyperparameter scripts, evaluation configurations, and
saved samples are kept in the project repository under
\code{training/}, \code{rewards/}, \code{verl\_bridge/}, and
\code{eval/}. The complete result table is in
\code{EVAL\_RESULT\_TABLE\_20260414.md}, and the
training-progress narrative is in \code{PROJECT\_REPORT.md}.
